Here is a complete proof. All indices are taken modulo \(n\).

**Notation.** For a compact convex set \(K\subset \mathbb{R}^2\) and a nonzero vector \(u\), write
\[
h_K(u):=\max_{x\in K}\langle u,x\rangle,
\qquad
F_K(u):=\{x\in K:\langle u,x\rangle=h_K(u)\}.
\]
Thus \(F_K(u)\) is the face of \(K\) exposed by \(u\).

---

### Basic lemma for Minkowski sums

For compact convex sets \(A,B\subset\mathbb{R}^2\) and \(u\neq 0\),
\[
F_{A+B}(u)=F_A(u)+F_B(u).
\]

**Proof.** First, if \(a\in F_A(u)\) and \(b\in F_B(u)\), then
\[
\langle u,a+b\rangle
=
h_A(u)+h_B(u)
=
h_{A+B}(u),
\]
so \(a+b\in F_{A+B}(u)\). Hence
\[
F_A(u)+F_B(u)\subseteq F_{A+B}(u).
\]

Conversely, let \(z\in F_{A+B}(u)\). Write \(z=a+b\) with \(a\in A\), \(b\in B\). Then
\[
h_A(u)+h_B(u)
=
h_{A+B}(u)
=
\langle u,z\rangle
=
\langle u,a\rangle+\langle u,b\rangle
\le h_A(u)+h_B(u).
\]
Therefore equality holds in both inequalities, so
\[
\langle u,a\rangle=h_A(u),
\qquad
\langle u,b\rangle=h_B(u),
\]
that is,
\[
a\in F_A(u),\qquad b\in F_B(u).
\]
Thus \(z\in F_A(u)+F_B(u)\), proving the reverse inclusion. \(\square\)

---

## Part \((1)\): the vertices of \(P+Q\)

Let
\[
x_i:=s_i+t_i.
\]

Fix \(i\), and let \(u\in C_i\). Since \(u\) lies in the open normal cone of \(P\) at \(s_i\),
\[
F_P(u)=\{s_i\}.
\]
By the defining property of \(t_i\) from Subproblem \(2\),
\[
F_Q(u)=\{t_i\}.
\]
Hence, by the lemma,
\[
F_{P+Q}(u)=F_P(u)+F_Q(u)=\{s_i+t_i\}=\{x_i\}.
\]
So \(x_i\) is an exposed point of \(P+Q\), hence a vertex of the polygon \(P+Q\).

Moreover, by Subproblems \(1\)–\(2\), the points
\[
\ell_i=s_i+t_i=x_i
\]
are the \(n\) support points arising in the extremal case, so they are distinct. Thus \(x_1,\dots,x_n\) are \(n\) distinct vertices of \(P+Q\).

Now we show there are no other vertices. Since \(P\) and \(Q\) are convex hulls of finite sets, \(P+Q\) is a convex polygon. Any vertex \(y\) of a polygon is exposed by all directions in a nonempty open arc of outer normals; call that arc \(D\). The normal fan of \(P\) consists of the \(n\) open cones \(C_1,\dots,C_n\) together with the \(n\) boundary rays corresponding to the edge normals of \(P\). Since \(D\) is open, it must meet some \(C_i\). For any \(u\in D\cap C_i\),
\[
\{y\}=F_{P+Q}(u)=\{x_i\},
\]
so \(y=x_i\).

Therefore the vertices of \(P+Q\) are exactly
\[
x_1,x_2,\dots,x_n.
\]

Finally, the cones \(C_1,C_2,\dots,C_n\) occur cyclically around the circle of directions, and for \(u\in C_i\) the exposed vertex is \(x_i\). Hence, as the outer normal rotates once around the circle, the vertices encountered on \(P+Q\) are exactly
\[
x_1,x_2,\dots,x_n
\]
in that cyclic order. This proves Part \((1)\).

---

## Part \((2)\): the face of \(Q\) at an edge normal of \(P\)

Fix \(i\), and set
\[
E_i:=[s_i,s_{i+1}],
\qquad
a_i:=s_{i+1}-s_i.
\]
Let \(u_i\) be any outer normal to the edge \(E_i\). Then
\[
F_P(u_i)=E_i,
\qquad
\langle u_i,a_i\rangle=0.
\]

Define
\[
G_i:=F_Q(u_i).
\]

### Step \((a)\): \(t_i\) and \(t_{i+1}\) both lie in \(G_i\)

Take a sequence \(u_m\in C_i\) with \(u_m\to u_i\). For each \(m\), the point \(t_i\) is the unique maximizer of \(\langle u_m,\cdot\rangle\) on \(Q\). Since \(Q=\operatorname{conv}(T)\), it is enough to compare on \(T\): for every \(t\in T\),
\[
\langle u_m,t_i\rangle\ge \langle u_m,t\rangle.
\]
Passing to the limit gives
\[
\langle u_i,t_i\rangle\ge \langle u_i,t\rangle
\qquad\text{for all }t\in T.
\]
A linear functional attains its maximum on \(\operatorname{conv}(T)\) at a point of \(T\), so this implies
\[
t_i\in F_Q(u_i)=G_i.
\]

The same argument, using a sequence from \(C_{i+1}\), shows that
\[
t_{i+1}\in G_i.
\]

So \(G_i\) contains both \(t_i\) and \(t_{i+1}\).

### Step \((b)\): \(G_i\) is exactly \([t_i,t_{i+1}]\)

Because \(G_i=F_Q(u_i)\), it lies in the support line
\[
L_i:=\{q\in\mathbb{R}^2:\langle u_i,q\rangle=h_Q(u_i)\}.
\]
This line is orthogonal to \(u_i\). Since \(E_i\) is also orthogonal to \(u_i\), the line \(L_i\) is parallel to \(E_i\), hence parallel to \(a_i\). Therefore \(G_i\) is a point or a segment contained in a line parallel to \(a_i\).

Now perturb \(u_i\) slightly in the \(\pm a_i\) directions. Since \(u_i\) is the common boundary direction between the cones \(C_i\) and \(C_{i+1}\), and \(a_i\) points from \(s_i\) to \(s_{i+1}\), for all sufficiently small \(\varepsilon>0\) we have
\[
u_i-\varepsilon a_i\in C_i,
\qquad
u_i+\varepsilon a_i\in C_{i+1}.
\]
Indeed, on the edge \(E_i\), the perturbation by \(-\varepsilon a_i\) favors \(s_i\), and the perturbation by \(+\varepsilon a_i\) favors \(s_{i+1}\); all other vertices of \(P\) stay strictly below the supporting line for small \(\varepsilon\).

Now take any \(q\in G_i\). Since \(\langle u_i,q\rangle=h_Q(u_i)\), we have
\[
\langle u_i-\varepsilon a_i,q\rangle
=
h_Q(u_i)-\varepsilon\langle a_i,q\rangle.
\]
Thus, on \(G_i\), maximizing \(\langle u_i-\varepsilon a_i,\cdot\rangle\) is equivalent to minimizing \(\langle a_i,\cdot\rangle\). But \(u_i-\varepsilon a_i\in C_i\), so its unique maximizer on \(Q\) is \(t_i\). Therefore \(t_i\) is the unique minimizer of \(\langle a_i,\cdot\rangle\) on \(G_i\).

Similarly,
\[
\langle u_i+\varepsilon a_i,q\rangle
=
h_Q(u_i)+\varepsilon\langle a_i,q\rangle,
\]
so, on \(G_i\), maximizing \(\langle u_i+\varepsilon a_i,\cdot\rangle\) is equivalent to maximizing \(\langle a_i,\cdot\rangle\). Since \(u_i+\varepsilon a_i\in C_{i+1}\), its unique maximizer on \(Q\) is \(t_{i+1}\). Therefore \(t_{i+1}\) is the unique maximizer of \(\langle a_i,\cdot\rangle\) on \(G_i\).

But \(G_i\) is a convex subset of a line parallel to \(a_i\). On such a set, the minimum and maximum of \(\langle a_i,\cdot\rangle\) occur exactly at the two endpoints. Hence the endpoints of \(G_i\) are precisely \(t_i\) and \(t_{i+1}\). Therefore
\[
G_i=[t_i,t_{i+1}].
\]

This proves the first assertion in Part \((2)\).

### Step \((c)\): the direction constraint

Since \(G_i=[t_i,t_{i+1}]\) lies on a line parallel to \(a_i\), there exists a real number \(\lambda_i\) such that
\[
t_{i+1}-t_i=\lambda_i a_i
=
\lambda_i\,(s_{i+1}-s_i).
\]
Because \(t_i\) is the \(a_i\)-minimum and \(t_{i+1}\) is the \(a_i\)-maximum on \(G_i\),
\[
\langle a_i,t_{i+1}-t_i\rangle\ge 0.
\]
Substituting \(t_{i+1}-t_i=\lambda_i a_i\) gives
\[
\lambda_i\langle a_i,a_i\rangle\ge 0.
\]
Since \(a_i\neq 0\), we have \(\langle a_i,a_i\rangle>0\), so
\[
\lambda_i\ge 0.
\]
Thus
\[
t_{i+1}-t_i
\]
is a nonnegative scalar multiple of
\[
s_{i+1}-s_i,
\]
with the degenerate case \(\lambda_i=0\) corresponding to \(t_i=t_{i+1}\).

This proves Part \((2)\).

---

## Stronger boundary identity for \(P+Q\)

At the same edge normal \(u_i\), the lemma gives
\[
F_{P+Q}(u_i)=F_P(u_i)+F_Q(u_i)
=
[s_i,s_{i+1}]+[t_i,t_{i+1}].
\]
Since
\[
t_{i+1}-t_i=\lambda_i a_i
\qquad\text{with}\qquad
\lambda_i\ge 0,
\]
we get
\[
[s_i,s_{i+1}]+[t_i,t_{i+1}]
=
\{s_i+t_i+(\alpha+\lambda_i\beta)a_i:0\le \alpha,\beta\le 1\}
=
[x_i,x_{i+1}].
\]
So the edge-normal direction of \(P\) exposes exactly the edge
\[
[x_i,x_{i+1}]
\]
of \(P+Q\).

---

Therefore:

\[
\boxed{\text{The vertices of }P+Q\text{ are exactly }x_i=s_i+t_i\text{ in cyclic order,}}
\]
and for each \(i\),
\[
\boxed{F_Q(u_i)=[t_i,t_{i+1}],\qquad
t_{i+1}-t_i=\lambda_i\,(s_{i+1}-s_i)\ \text{for some }\lambda_i\ge 0.}
\]

